% !TEX root = ../main.tex
% !TEX TS-program = pdflatex
% !TEX encoding = UTF-8
% !TEX spellcheck = en_US

%%%% ---- LaTeX File ---- %%%%
\providecommand{\PathP}{./parts/} % Path to  Resources folders
\providecommand{\PathR}{./resources/} % Path to  Resources folders
\providecommand{\PathSO}{./sources/} % Path to Sources folders
\providecommand{\PathS}{./supplies/} % Path to Supplies folders
\providecommand{\lang}{english}	% english, french
\providecommand{\BuildTrigger}{buildnew} % Trigger to build or not standalone parts.

\makeatletter
	\def\input@path{ {\PathSO}{./another-directory/} } % to be able to find package in source folders.
	\ifx\c@levelstanda\undefined % to create a new-counter if it doesn't exists.
	    \newcounter{levelstanda}
	\else
	    {}% it already exist
	\fi
\makeatother


%%%%---- DOCUMENT-CLASS -------------------------%%%%
\documentclass[a4paper]{article}  % or another class
	\usepackage[l2tabu, orthodox]{nag}	% Helps to correct the script
 
%%%%---- OS & FONT -------------------------%%%%
	\usepackage[utf8]{inputenc} % Input keyboard encoding: latin1=windows, utf8=linux, applemac=Mac
	\usepackage[T1]{fontenc} 	% Output font encoding: T1 for Europeans languages.

%%%%---- PERSONAL Custom ----------------------%%%%
    \newcommand{\Pathbasebib}{\PathR/ref_RN.bib,\PathR/ref_XX.bib}% Careful no space!
    \newcommand{\PathSymbolsEntries}{\PathR/UsersSymbolsAcronyms.bib}
	\input{\PathR/UsersPackagesDefs.tex}
	\input{\PathR/UsersMathsCommands.tex}
        \geometry{margin=2.5cm}

    
%%%%%---- JOURNAL SPECS -------------------------%%%%	
    %\journal{Journal of XX}
    \title{My eXAMPLE Template}
    \author{Romain \famName{Noël}}
    \date{\today}


%%%%%%%%%%%%%%%  MAIN DOCUMENT  %%%%%%%%%%%%%%%%%%%%%
\begin{document} 
\stepcounter{levelstanda}
    
%% FRONTMATTER
\ifnum\value{levelstanda}>1
	{}% Not in the main.
\else
    \maketitle
\fi


%% MAINMATTER
\section{Introduction}
    \label{sec:intro}  % \label{} allows reference to this section
    Example of an acronym managed by 'glossaries' package: \gls{lbm} from \cite{Bhatnagar_1954,Chapman_1970}.
    But the second call of \gls{lbm} appears shorter.

    Example of an equation using 'derivative' package and glossaries (for symbols):
    \begin{equation}
    	\label{eq:boltzmann}
    	\pdv{\gls{BoltzDist}}{\gls{time}}
    	+\gls{parti_spe}\cdot\gradN[\gls{posi}]{\gls{BoltzDist}} + \gls{forcg}\cdot\gradN[\gls{parti_spe}]{\gls{BoltzDist}} =\CollOp{\gls{BoltzDist},\gls{BoltzDist}}
        \ .%
    \end{equation}

     
\section{Sate of the Art}
\subsection{This is a subsection}

    \begin{figure}[htb]
        \centering
        \begin{subfigure}[b]{0.475\textwidth}
            \centering
            \includegraphics[width=\linewidth, draft]{\PathS/example/D1Q5.pdf}
            \caption{Something draft a.}
            \label{fig:draftA}
        \end{subfigure}
        \begin{subfigure}[b]{0.475\textwidth}
            \centering
            \includegraphics[width=\linewidth, draft]{\PathS/example/D1Q5.pdf}
            \caption{Something else draft b.}
            \label{fig:draftB}
        \end{subfigure}
        \caption{Example of subfigures where (\subref{fig:draftA}) is a first element and (\subref{fig:draftB}) is a second one.}
        \label{fig:subfig}
    \end{figure}

    \begin{figure}[htbp]
        \centering
		% \UpdateLineAndBaseLength{} % set the local
		% \UpdatePlotSize{width=0.6\string\linewidth}{}
        % \begin{nomemoize}
		% \includestandalone[mode=\BuildTrigger, %build, %
		% 		% width=0.99\linewidth
		% 	]{\PathS/example/plot_test}%.tex
        % \end{nomemoize}
		% \renewcommand{\plotsize}{}
	    \input{\PathS/example/plot_test}%.tex
        \caption{Test of simple plot using TikZ+CSV+glossaries.}
		\label{fig:test}
	\end{figure}
    
	
	\begin{figure}[htbp]
		\centering
		% \includestandalone[mode=\BuildTrigger, %build,
		% 		%width=0.99\linewidth
		% 	]{\PathS/example/D1Q5-stda}%.tex
	    \input{\PathS/example/D1Q5-stda}%.tex
    	\caption{Test of a simple sketch using Inkscape+\LaTeX+glossaries}
		\label{ch1:fig:D1Q5}
	\end{figure}

    \begin{table}[htbp]
		% %\UpdateLineAndBaseLength{} % set the local
		% %\UpdatePlotSize{width=0.6\string\linewidth}{}
		% % \centering
		% \includestandalone[mode=\BuildTrigger, %build, %
		% 		%width=0.99\linewidth
		% 	]{\PathS/example/table_csv_test}%.tex
		% %\renewcommand{\plotsize}{}
		\input{\PathS/example/table_csv_test}%.tex
        \caption{Test of table tabularray+CSV+glossaries+SIunitX.}
		\label{tab:test}
	\end{table}

    \begin{table}[htbp]
		% %\UpdateLineAndBaseLength{} % set the local
		% %\UpdatePlotSize{width=0.6\string\linewidth}{}
		% % \centering
		% \includestandalone[mode=\BuildTrigger, %build, %
		% 		%width=0.99\linewidth
		% 	]{\PathS/example/table_simple-stda}%.tex
		% %\renewcommand{\plotsize}{}
		\input{\PathS/example/table_simple-stda}%.tex
        \caption{Test of table standalone manual tabularray.}
		\label{tab:test_manual}
	\end{table}
    
	\begin{algorithm}{Example using 'pseudo': Euclid's algorithm, \pr{Euclid}(a, b)}{euclid}
    %     \includestandalone[mode=\BuildTrigger, %build,
    %         %width=0.99\linewidth
    %     ]{\PathS/example/algo}%.tex
    %     % \input{\PathS/example/algo}
        \input{\PathS/example/algo}%.tex
        % \caption{Example using 'pseudo': Euclid's algorithm, \pr{Euclid}(a, b)}
        % \label{alg:euclid}
    \end{algorithm}

    And here is a ref to Euclid in \cref{alg:euclid}.
	
	\FloatBarrier
	
    \begin{lstlisting}[language=C++, frame=single, caption=An extract of \lstinline$C++$ code using 'listing'., gobble=8]
        #include <stdio.h>
        #define N 10
        /* Block
         * comment 
         */
        
        int main() {
            int i;
        
            // Line comment.
            puts("Hello world!");
            
            for (i = 0; i < N; i++)
            {
                puts("LaTeX is also great for programmers!");
            }
        
            return 0;
        }
    \end{lstlisting}

    \lstinputlisting[language=Python, frame=single, 
        caption=Another example including directly the \lstinline$Python source$ file.,
        firstline=2,lastline=7
    ]{\PathS/example/file.py}

    Otherwise, you can also use lstinline to refer to \lstinline{some code}.
    
    Something after.
    \chhighlight[id=RN]{highlighted}.

    You can also hide some text using comment (and add a post-it to remember that it is here):
    \begin{comment}this will not appear\end{comment}
    \chcomment[id=RN]{Here there is some hidden latex code.}


\section*{Acknowledgments}       
 
    This unnumbered section is used to identify those who have aided the authors in understanding or accomplishing the work presented and to acknowledge sources of funding.  
    
    
%% BACK-MATTER
%\appendix
\ifnum\value{levelstanda}>1
	{}% Not in the main.
\else
    % References    
    \bibliography{\Pathbasebib} % for bibtex or natbib
    %\printbibliography     % for biblatex
    
	%%% To comment for the final version
	%\tableofcontents
	\printunsrtglossary[type=acronym]
	\printunsrtglossary[type=symbol]
\fi

\addtocounter{levelstanda}{-1}
\end{document}
%EoF